\section{Literature Review}
\label{sec:literature}

This review does not attempt a comprehensive survey of electricity price
forecasting, for which recent reviews are available
\cite{lago2021forecasting,weron2014electricity}. It focuses instead on the
question that motivates this thesis: how do existing day-ahead price models
represent their inputs, and is that representation admissible in an
operational, pre-gate-closure setting? Table~\ref{tab:selected_literature} is
the central output of the review. Adapting the feature overview of
\textcite{eiser2026operational}, it records for each study which input features
enter the price model, whether direct weather is used, whether all inputs are
observable before gate closure, and whether the study is open and reproducible,
and contrasts these with this thesis.

\noindent\textbf{Day-Ahead Price Forecasting.}
\textcite{weron2014electricity} provides a broad review of electricity price
forecasting and emphasizes the volatility, seasonality, and short-run
dependence of spot prices. \textcite{lago2021forecasting} extend this view by
reviewing modern algorithms, proposing an open benchmark for day-ahead
electricity price forecasting, and showing that strong empirical designs
usually combine price lags, calendar variables, and exogenous system
information.
\textcite{kath2018quarterHourly} study quarter-hourly electricity price
forecasts and show that more accurate forecasts at this resolution can create
economic value. The quarter-hourly resolution of the \mbox{DE-LU} price target
follows from the market design rather than from a modeling choice. Together,
these contributions motivate the use of lagged prices, calendar variables, and
simple naive price benchmarks in this thesis.

\noindent\textbf{Load and Renewable Generation in Price Models.}
Load, solar, and wind generation are fundamental drivers of the
price, because they describe the physical supply-and-demand state of the power
system \cite{do2021residualDemand,sensfuss2008meritOrder}.
\textcite{sensfuss2008meritOrder} document the merit-order effect for Germany
and show how renewable feed-in can reduce spot prices by shifting more
expensive conventional generation out of the market. \textcite{pape2016fundamentals}
use a fundamental supply-and-demand model of the German day-ahead and intraday
market to examine how far load, renewable generation, and the conventional
supply stack explain observed price variations, while \textcite{trebbien2023merit} use
explainable machine learning to show that German electricity prices depend on
interacting drivers beyond a simple merit-order view. Most directly,
\textcite{maciejowska2021enhancing} study the role of load, wind, and solar
generation information in day-ahead electricity price forecasting and show that
including these forecast variables, and improving their accuracy, yields more
accurate price forecasts. Their comparison is between price models with and
without these forecasts, rather than against a model that already uses direct
weather variables. Whether self-generated load, solar, and wind forecasts still
improve a price model that already includes raw weather therefore remains open and is tested in \hyperlink{rq:price_impact_generated_forecasts}{RQ2}.

A separate line of research brings weather variables directly into the
price model. \textcite{ludwig2015bigdata} forecast German spot prices using
historical prices and a large set of weather variables, while
\textcite{sgarlato2023weatherPredictions} show that meteorological forecasts
can improve German electricity price forecasts beyond the day-ahead horizon.
These papers show that direct weather-to-price representations are feasible.
They also make the representation question explicit: weather can enter a price
model directly, or it can enter indirectly through load, solar, and wind
forecasts. Using forecasts as this indirect representation raises two questions
that this thesis addresses. First, the forecasts must be accurate, and
\hyperlink{rq:component_forecast_quality}{RQ1} benchmarks the OBTF
load, solar, and wind forecasts against the \mbox{ENTSO-E} forecasts. Second,
they must capture enough of the weather-driven state to stand in for raw
weather, which \hyperlink{rq:price_impact_generated_forecasts}{RQ2} examines.

\noindent\textbf{Forecasting Load, Solar, and Wind Generation.}
Generating these forecasts is itself a forecasting task with an established
literature for each target. For load,
\textcite{hong2016probabilisticLoad} show why calendar effects, recent load
observations, and weather variables are central predictors. For photovoltaic
generation, \textcite{antonanzas2016pvForecasting} identify radiation, cloud
cover, temperature, and plant characteristics as key inputs, while
\textcite{foley2012windForecasting} review wind power forecasting methods and
emphasize the role of wind speed, wind direction, atmospheric conditions, and
turbine locations.

\noindent\textbf{Public Forecasts and Operational Timing.}
Many \mbox{DE-LU} price forecasting models rely on system forecasts published on the
\mbox{ENTSO-E} Transparency Platform. In their review,
\textcite{eiser2026operational} find that these models commonly use the
day-ahead load forecast and, in several cases, the day-ahead solar and wind
forecasts, while direct weather variables are not listed as a separate input
category. Table~\ref{tab:selected_literature} compares the day-ahead price-forecasting
studies most directly related to this thesis. It places these models, which use the published
\mbox{ENTSO-E} load and renewable-generation forecasts as price-model inputs,
next to the weather-direct studies and the operational
pipeline of \textcite{eiser2026operational}, and records their input features,
use of direct weather, pre-gate-closure admissibility, and openness. For price modeling, these published load and renewable-generation forecasts act
as summary variables that compress weather, calendar, and system information into
expected load, solar, and wind generation, so a model that uses them
and omits
direct weather treats them as the relevant compression of weather information.

In hindsight, every published forecast is available in the data, including
forecasts that were not yet observable when the price forecast would have been
made. A valid backtest must therefore apply the same restriction as an
operational forecast, namely that every input is observable before the 12:00
gate closure \cite{epexSpotPowerMarket}. Under the transparency regulation the public load
forecast is published no later than two hours before gate closure and is
observed around 10:30, so it is admissible before gate closure
\cite{euRegulation543_2013,entsoeTransparency}. The forecast values are published, but the models that produce them are not
disclosed. Their errors can be measured against the realized values, but an error level
alone does not show how good a forecast is compared with alternatives. \hyperlink{rq:component_forecast_quality}{RQ1} therefore benchmarks
these forecasts against forecasts built from open data. The day-ahead price model of
\textcite{uniejewski2025smoothing}, for instance, relies on price lags, calendar
variables, and this admissible load forecast, without any
renewable-generation forecasts.
The \mbox{ENTSO-E} solar and wind day-ahead forecasts, by contrast, are not
publicly observable before gate closure, with a publication requirement no later
than 18:00 on day \(d-1\)
\cite{kleinebrahm2026energyarena,euRegulation543_2013}, so this thesis uses them
only as benchmarks. Faced with this constraint,
\textcite{eiser2026operational} substitute spatially resolved weather for the
unavailable renewable-generation forecasts rather than forecasting renewable generation
explicitly, and the Energy-Arena benchmark \cite{kleinebrahm2026energyarena}
enforces the same pre-gate-closure timing across contributors. The reviewed
studies do not systematically examine the trade-off between forecasting earlier
and using later-arriving information before gate closure.
This thesis addresses both gaps. In
\hyperlink{rq:component_forecast_quality}{RQ1} it generates its own OBTF load, solar, and wind forecasts from information available before gate closure and benchmarks
them against the corresponding \mbox{ENTSO-E} forecasts. In
\hyperlink{rq:cutoff_sensitivity}{RQ3} it then examines how price forecast
accuracy changes as the forecast creation time moves toward gate closure.

\noindent\textbf{Weather and Spatial Information.}
Weather forecasts are the main input for load, solar, and wind forecasts
\cite{hong2016probabilisticLoad,antonanzas2016pvForecasting,foley2012windForecasting}.
The Icosahedral Nonhydrostatic (ICON) model framework is described by
\textcite{zaengl2015icon}, and the Deutscher Wetterdienst (DWD) provides operational weather
forecast products as open data \cite{dwdIconForecastData}. Open-Meteo is used as an
additional open weather application programming interface (API) in the pipeline \cite{openMeteoForecastApi}. The
weather fields have to be weighted by where demand and generation are located.
Population data from Eurostat therefore weight the weather for the OBTF load model,
and installed renewable capacity from the Marktstammdatenregister (MaStR) weights it for
the OBTF solar and wind models \cite{eurostatGisco2021,bundesnetzagenturMastr}.

\begin{table}[htbp]
  \centering
  \footnotesize
  \setlength{\tabcolsep}{4pt}
  \renewcommand{\arraystretch}{1.2}
  \caption[Input-feature representation, admissibility, and openness across studies.]
  {Input features, pre-gate-closure admissibility, and openness of selected German and
  \mbox{DE-LU} day-ahead price-forecasting studies compared with this thesis. A check mark denotes a
  reported property, a dash an absent or unreported one. The \emph{ENTSO-E} and
  \emph{Self-generated} columns mark the use of published \mbox{ENTSO-E} forecasts or of
  forecasts produced from admissible data. \emph{Pre-GC}: all inputs observable before the 12:00
  gate closure. \emph{Data}, \emph{Code}: released dataset or pipeline and public code.
  \emph{Benchmark}: results reported on Energy-Arena. RES: renewable energy sources.
  EUA: EU emission allowances.
  \textsuperscript{a}States a 12:00 cutoff but uses the post-closure \mbox{ENTSO-E} RES forecast.
  \textsuperscript{b}Bias-corrects the published \mbox{ENTSO-E} forecasts.
  \textsuperscript{c}Correction of the published \mbox{ENTSO-E} load forecast.}
  \label{tab:selected_literature}
  \begin{tabular}{@{}l *{12}{c}@{}}
    \toprule
    & \multicolumn{4}{c}{General inputs} & \multicolumn{4}{c}{Load / RES forecast}
    & \multicolumn{4}{c}{Operational \& openness} \\
    \cmidrule(lr){2-5}\cmidrule(lr){6-9}\cmidrule(l){10-13}
    Study
    & \rotatebox{90}{Price lags}
    & \rotatebox{90}{Direct weather}
    & \rotatebox{90}{Fuel/EUA}
    & \rotatebox{90}{Calendar}
    & \rotatebox{90}{ENTSO-E load}
    & \rotatebox{90}{ENTSO-E RES}
    & \rotatebox{90}{\makecell[l]{Self-generated\\load}}
    & \rotatebox{90}{\makecell[l]{Self-generated\\RES}}
    & \rotatebox{90}{Pre-GC}
    & \rotatebox{90}{Data}
    & \rotatebox{90}{Code}
    & \rotatebox{90}{\makecell[l]{Benchmark\\(Energy-Arena.org)}} \\
    \midrule
    \textcite{trebbien2023probabilistic}
      & \checkmark & -- & \checkmark & --
      & \checkmark & \checkmark & -- & --
      & --\textsuperscript{a} & -- & -- & -- \\
    \textcite{marcjasz2023distributional}
      & \checkmark & -- & \checkmark & \checkmark
      & \checkmark & \checkmark & -- & --
      & -- & \checkmark & \checkmark & -- \\
    \textcite{hilger2024multivariate}
      & \checkmark & -- & -- & --
      & \checkmark & \checkmark & -- & --
      & -- & -- & -- & -- \\
    \textcite{maciejowska2021enhancing}
      & \checkmark & -- & -- & \checkmark
      & \checkmark & \checkmark\textsuperscript{b} & -- & --
      & -- & -- & -- & -- \\
    \textcite{uniejewski2025smoothing}
      & \checkmark & -- & -- & \checkmark
      & \checkmark & -- & -- & --
      & \checkmark & -- & -- & -- \\
    \textcite{ludwig2015bigdata}
      & \checkmark & \checkmark & -- & --
      & -- & -- & -- & --
      & \checkmark & -- & -- & -- \\
    \textcite{sgarlato2023weatherPredictions}
      & \checkmark & \checkmark & -- & --
      & -- & -- & -- & --
      & \checkmark & -- & -- & -- \\
    \textcite{eiser2026operational}
      & \checkmark & \checkmark & -- & \checkmark
      & \checkmark & -- & -- & --
      & \checkmark & \checkmark & \checkmark & \checkmark \\
    \textbf{This thesis}
      & \checkmark & \checkmark & -- & \checkmark
      & -- & -- & \checkmark\textsuperscript{c} & \checkmark
      & \checkmark & \checkmark & \checkmark & \checkmark \\
    \bottomrule
  \end{tabular}
\end{table}

\FloatBarrier

\noindent\textbf{Operational Benchmarking.}
The operational requirements emphasized in this review are brought together in
the Energy-Arena operational forecasting benchmark
\cite{kleinebrahm2026energyarena}. Energy-Arena hosts day-ahead forecasting
challenges for the \mbox{DE-LU} and Austrian markets, with separate load, solar, wind, and price challenges as well as challenges
tied to different forecast creation times. It scores each submission against
those of the other contributors on the realized market outcome under a common
gate-closure deadline, so that all participants face the same operational
information constraint and their forecasts can be compared directly on a public
leaderboard. It thereby turns the pre-gate-closure
setting studied here into an external, leakage-controlled testbed. This thesis
submits its OBTF load, solar, and wind forecasts and its resulting
price forecasts to the benchmark at different cutoff times, where
they remain publicly accessible and
reproducible and can be reused by others, for example as members of forecast
ensembles. Energy-Arena is thus where the elements reviewed above come together
in practice: transparent open-data inputs, forecasts produced under an
admissible information set, a fixed pre-gate-closure timing, and open, reusable
outputs.

\noindent\textbf{Summary and Positioning.}
That load, solar and wind generation, weather, calendar variables, and
price lags matter for day-ahead price forecasting is well established. Three less
obvious points motivate this thesis. First, the representation of load and
renewable generation matters, not only their inclusion. \textcite{maciejowska2021enhancing}
show that how load and renewable generation are predicted, and how accurate those
predictions are, changes the price forecast. Second, how weather should be represented for a price model is not
systematically studied. Weather usually enters only indirectly, through
published load and renewable-generation forecasts, so the model inherits
whatever those forecasts capture, and the studies that do use weather directly
neither compare it against that indirect representation nor test whether the
load and renewable-generation forecasts add value on top of direct weather. This thesis
takes up that representation question, comparing a weather-based price model
with one that adds the OBTF load, solar, and wind forecasts and testing, through the
weather-free variant, whether those forecasts can substitute for raw weather
(\hyperlink{rq:price_impact_generated_forecasts}{RQ2}).

Third, and most importantly, the operational information set is often left
underspecified. A cutoff time alone does not define the experimental setup. What
matters is which inputs are actually observable at that time, yet the reviewed
studies often do not establish which inputs are admissible by then. Some
therefore include inputs that are published only after
gate closure, such as the \mbox{ENTSO-E} renewable-generation forecasts
\cite{trebbien2023probabilistic,marcjasz2023distributional,hilger2024multivariate},
while informative and clearly admissible signals, such as early exchange prices
\cite{ziel2015exaa}, are left out without comment. The first opens the setup to information leakage,
the second discards admissible information, and both follow from not
establishing which inputs are admissible. This
thesis instead defines the admissible information set explicitly and tracks how
it grows through the morning up to gate closure, so that every comparison uses
only information that was genuinely available at the forecast creation time.

None of the reviewed operational \mbox{DE-LU} studies generates dedicated load,
solar, and wind forecasts from admissible information and uses them as
price-model inputs, and none systematically examines how price forecast quality
changes with the forecast creation time before gate closure. This thesis
addresses both gaps, and does so transparently: it uses open data,
releases its full pipeline as open source, and submits its forecasts to the
Energy-Arena. These properties are contrasted across
the reviewed studies in Table~\ref{tab:selected_literature}.

This gap leads directly to the methodology. The thesis generates OBTF load, solar, and wind forecasts from
admissible information, compares them with the corresponding \mbox{ENTSO-E}
day-ahead forecasts, and tests whether they add value to the day-ahead price
model. Finally, it evaluates how price forecast performance
changes as the forecast creation time moves from 07:00 toward 12:00.
